\answer{ex:03:1}{$g_E=g_L$에서 $35\to17.5\mV$; $g_E=4g_L$에서 $56\to40\mV$이고, 뺄셈이라면 $38.5\mV$일 것이다. 분로는 $g_E$를 $3$으로 나눈다. 창은 절반으로 좁아진다($\tau_{\text{eff}}$가 $2$배 줄어든다).}
\answer{ex:03:2}{$\operatorname{tr}=\frac{w_{EE}-1}{\tau_E}-\frac1{\tau_I}$, $\det=\frac{1-w_{EE}+w_{EI}w_{IE}}{\tau_E\tau_I}$; 경계에서 $\omega=\sqrt{\det}$, $\tau_I=20\ms$, 주기 $73\ms$.}
\answer{ex:03:3}{$d_c=\tau_E\arccos(-1/G)/\sqrt{G^2-1}$, 주기 $2\pi\tau_E/\sqrt{G^2-1}\in(2d_c,4d_c)$. $G=2$: $d_c=12.1\ms$, 주기 $36\ms$; $G=5$: $d_c=3.6\ms$, 주기 $12.8\ms$. 딱 빠른 리듬의 대역이다.}
\answer{ex:03:4}{(a)~올릴 때는 $I_2=\beta I_1=2$에서, 내릴 때는 $I_2=I_1/\beta=0.5$에서 전환된다. 폭 $1.5$의 고리이다. (b)~$\varepsilon$이 한 자릿수 줄 때마다 $\tau\ln10/(\beta-1)=2.3\tau$씩 늘어난다.}
\answer{ex:03:5}{중개자 세 개, $d_k=5$, $10$, $15\ms$: 금지가 $15\ms$를 덮어야 하고, 하나가 $5\ms$씩 덮는다.}
\answer{ex:03:6}{뉴런 $n+M+1=3+4+1=8$개. 두 개로는: \nt{GABA} $=[x+y+z\ge2]$, 출력 $=[x+y+z-2\,\nt{GABA}\ge1]$.}
\answer{ex:03:7}{$\binom84=70<100\le\binom94=126$: 중개자 $9$개(슈페르너 정리). 직접 연결이 없으면 $100$개: $\beta(J-I)$의 계수가 $n$이다.}
\answer{ex:03:8}{$h_c=\frac1{2w_{EE}}+\frac{w_{EI}w_{IE}-w_{EE}}{4w_{EE}^2}=\frac7{18}\approx0.39$; 시뮬레이션에서는 $h=0.38$과 $0.40$ 사이에서 부호가 바뀐다.}
